\section{Trace-one transfer of an existing ordinary cusp form}
\label{sec:cubic-tame-cusp-transfer}

This is a supporting result, not a construction of the missing
cubic Type I seed.  Fix the imaginary quadratic Picard setting of
AC-G67-boundary, a split coefficient prime $p>3$, and an auxiliary
split prime $q\ne p$ with $p\mid q-1$.  Retain a sufficiently fine
prime-to-$p$ level away from $q$ so that the open moduli schemes
are fine.  At $q$ write $I^-$ for the opposite Iwahori and $I_1^-$
for the kernel of its diagonal reduction.  The similitude factor
is kept in $1+q\mathbb Z_q$.  Thus
\[
 A=I^-/I_1^-=(\mathbb F_q^\times)^3=C\times A',
 \qquad C=A[p^\infty],\qquad p\nmid |A'|.
                                                        \tag{CT.1}
\]
Let $\mathcal O$ be a coefficient discrete valuation ring,
with uniformizer $\varpi$, containing the values of a character
$\lambda:A\to\mathcal O^\times$ of $p$-power order.  In
particular $\lambda|_{A'}=1$ and $\lambda\equiv1\pmod\varpi$.

\begin{proposition}[AC-G71: integral cusp transfer at unchanged weight]
\label{AC-G71}
Let $\kappa$ be a weight for which the actual ordinary classical
control isomorphism holds at the two tame levels under
consideration.  For example, one can use the classical weight
set in Hsieh (4.3), with Theorem~4.19, Theorem~4.21 and
Corollary~4.23 at these fixed levels.  Let $f$ be an integral
ordinary classical cusp form of weight $\kappa$, invariant
under $A$ at level $I_1^-$, whose reduction is nonzero.  Then
there is an integral ordinary classical cusp form $f_\lambda$
of the same weight at that level such that
\[
 a f_\lambda=\lambda(a)f_\lambda\quad(a\in A),
                 \qquad f_\lambda\equiv f\pmod\varpi.
                                                        \tag{CT.2}
\]
The congruence takes place in the same integral ambient
ordinary space.  If $f$ is a simultaneous eigenform for the
commuting integral Hecke operators away from the changed tame
level, together with the ordinary $p$-operators, a classical
ordinary cuspidal eigensystem in the exact $\lambda$-space is
congruent to that of $f$, after finite coefficient extension.

For
\[
       \lambda(t_1,t_2,t_3)=\vartheta_q(t_1)^{-1}
                                      \vartheta_q(t_3),   \tag{CT.3}
\]
where $\vartheta_q$ has order $p^r>1$, the resulting local
automorphic constituents have the finite principal-series
inertial type of AC-G67-type and zero monodromy.  A statement
about attached Galois inertia uses the local--global
compatibility input retained there.

Neither existence of the starting $f$ in the required cubic
branch nor global irreducibility of the resulting rank-three
Galois representation is an assertion of this proposition.
The classical-control hypothesis is at the unchanged weight;
it is not replaced by density of some other set of weights.
\end{proposition}
\begin{proof}
\proofstep{The deck action is free on the whole ordinary minimal space.}
On the open fine moduli scheme, changing the prime-to-$p$
$q$-level structure is a finite etale level cover: the
$q$-torsion of an abelian scheme is finite etale because $q$
is invertible, and a full level structure trivializes this
finite etale module locally in the etale topology.  The
intermediate $I_1^-$ to $I^-$ cover is therefore the torsor
for the finite constant group $A$.  Fineness eliminates
automorphisms of the underlying moduli object which could
identify two such level structures.  In particular $C$ acts
freely on the open moduli locus, in every characteristic
under consideration.

On the geometric minimal boundary, AC-G67-boundary supplies
the separate stabilizer calculation.  Its rational Levi
eigenvalues must be units at every finite place of the
imaginary quadratic field, hence roots of unity of order
dividing $4$ or $6$.  Their $p$-primary diagonal reductions
are trivial.  Thus a fixing element of $C$ is the identity.
The same calculation applies in characteristic $p$ through
the finite etale zero-dimensional boundary models, as proved
there.  These two loci exhaust the minimal space, so the
$C$-action is free on it and on its ordinary open.

Choose an $A$-invariant affine ordinary open by taking the
product of the translates of a lifted Hasse invariant.  Its
$p$-adic completion is the usual ordinary formal locus.
This is the affine minimal ordinary space used in the proof
of Hsieh Theorem~4.26.  A free action of the finite constant
group $C$ gives a finite etale $C$-torsor onto its quotient;
one may apply Stacks, Tags \texttt{07S7} and \texttt{0B87}.
The quotient of the affine open is affine.  The same
conclusions hold after coefficient extension and $p$-adic
completion.

\proofstep{Choose a global integral trace-one function.}
Write $B$ for the ring of functions on this ordinary formal
cover and $B_0$ for the quotient ring.  Etale locally the
cover is a disjoint union indexed by $C$, and
\[
                   \operatorname{Tr}:B\longrightarrow B_0
\]
is the sum of the coordinates.  It is consequently
surjective as a map of quasi-coherent sheaves.  Affineness
gives a global $a\in B$ satisfying
\[
                       \sum_{c\in C}c(a)=1.              \tag{CT.4}
\]
This also follows before completion and then persists after
completion.  Average $a$ over $A'$ using $|A'|^{-1}$; because
$A$ is abelian, (CT.4) is preserved and $a$ becomes
$A'$-invariant.  Its pullback to the ordinary Igusa tower has
weight zero.

\proofstep{Transfer the cusp form.}
On integral $p$-adic modular forms put
\[
 P_a(m)=\sum_{c\in C}\lambda(c)^{-1}c(a m),
                        \qquad f_\lambda=e_{\rm ord}P_a(f).
                                                        \tag{CT.5}
\]
The character calculation uses only a change of summation
variable and gives $dP_a(m)=\lambda(d)P_a(m)$ for $d\in C$.
For $m=f$, the $A'$-invariance of $a$ and $f$ gives the
remaining part of the $A$-character.  The ordinary projector
commutes with these tame deck operators.

Multiplication by $a$ preserves weight, since $a$ has weight
zero.  It also preserves cuspidality: the cusp sheaf is the
kernel of restriction to the minimal boundary, and
$(a f)|_\partial=a|_\partial f|_\partial=0$.
Every deck operator preserves the boundary, and the
ordinary projector commutes with restriction.  Hence (CT.5)
is an ordinary $p$-adic cusp form of weight $\kappa$.
It is integral.  Moreover
\[
 P_a(f)\equiv\sum_{c\in C}c(a)c(f)
                =\left(\sum_{c\in C}c(a)\right)f=f
                       \pmod\varpi.                     \tag{CT.6}
\]
The projector commutes with reduction and fixes $f$, so
(CT.2) follows.  The stipulated control isomorphism at
$\kappa$ identifies this ordinary $p$-adic form with a
classical one.  No change of weight occurs in this step.

\proofstep{Lift the congruent eigensystem in the exact character lattice.}
Let $M$ be the finite free integral ordinary classical cusp
module at the fine level and fixed weight, and let
\[
             M_\lambda=\{m\in M:a m=\lambda(a)m
                                             \ (a\in A)\}.
\]
It is finite free over the discrete valuation ring and is
$\varpi$-saturated in $M$: if $\varpi m\in M_\lambda$,
then $\varpi(a-\lambda(a))m=0$, hence
$(a-\lambda(a))m=0$, since $M$ is torsion-free.  Thus
$M_\lambda/\varpi M_\lambda$ injects into $M/\varpi M$.
By (CT.2), the image of $f_\lambda$ is the nonzero
eigenvector $\bar f$.  The same saturation shows that the
Hecke eigenrelations hold in $M_\lambda/\varpi M_\lambda$,
not merely in the larger reduction.

The algebra $\mathbb T$ generated by the chosen commuting
integral operators in $\operatorname{End}_{\mathcal O}
(M_\lambda)$ is finite and torsion-free over $\mathcal O$.
Localize at the maximal ideal specified by $\bar f$.
Its nonzero finite torsion-free module has a nonzero generic
fiber.  A maximal ideal of that generic fiber, after a
finite coefficient extension, supplies a simultaneous
eigencharacter specializing to the given maximal ideal.
Equivalently this is the finite-module Deligne--Serre
lifting argument.  It occurs in the classical ordinary
cuspidal module $M_\lambda$, proving the eigensystem claim.
Finally apply the exact local character-type conclusion
of AC-G67-type to (CT.3).
\end{proof}

\begin{remark}[The local character does not eliminate a two-plus-one parameter]
At the split prime $q$, the type (CT.3) admits the reducible
Weil--Deligne parameter
\[
   \bigl((\vartheta_q^{-1}\alpha)\oplus
                (\vartheta_q\beta)\bigr)\ \oplus\ \zeta,
                   \qquad N=0,                          \tag{CT.7}
\]
where $\alpha,\beta,\zeta$ are unramified characters.
Its two-dimensional block and its one-dimensional block
have precisely the same three distinct inertia characters
as the asserted rank-three principal-series type.  Thus
that type, its zero monodromy, and the integral congruence
do not themselves rule out an endoscopic $2+1$ constituent.
The analogous trace-one argument for an existing ordinary
modular cusp form can impose the pair
$(\vartheta_q^{-1},\vartheta_q)$ on its rank-two block;
an assertion about its transfer to a particular global
unitary packet still requires that packet's actual
multiplicity and local-label inputs.  No global packet
existence is deduced from the local display (CT.7).
\end{remark}
